% ECONOMICS_PROBLEM: {"id": "C78-293", "title": "Growing numbers of changed agents: correcting an overbroad complexity promise", "jel_code": "C78", "source_id": "OP-0199"}
% !TeX program = xelatex
\documentclass[11pt,a4paper]{article}
\usepackage[margin=23mm]{geometry}
\usepackage{fontspec}
\setmainfont{Latin Modern Roman}
\usepackage{amsmath,amssymb,mathtools}
\usepackage{xcolor,enumitem,booktabs,longtable,array,xurl,hyperref}
\definecolor{muted}{HTML}{53636C}
\hypersetup{colorlinks=true,urlcolor=blue,linkcolor=blue}
\setlength{\parindent}{0pt}
\setlength{\parskip}{5pt}
\newcommand{\R}{\mathbb R}
\newcommand{\E}{\mathbb E}
\newcommand{\N}{\mathbb N}
\newcommand{\ind}{\mathbf 1}
\newcommand{\Var}{\operatorname{Var}}
\newcommand{\Cov}{\operatorname{Cov}}
\newcommand{\supp}{\operatorname{supp}}
\DeclareMathOperator*{\argmin}{arg\,min}
\DeclareMathOperator*{\argmax}{arg\,max}
\newcommand{\entrymeta}[3]{{\sffamily\small\color{muted}\textbf{\texttt{#1}}\quad|\quad #2\par #3\par}}
\newcommand{\Problemref}[1]{\texttt{#1}}
\newcommand{\JELref}[2]{\texttt{#2} (JEL \texttt{#1})}
\begin{document}
\section*{Growing numbers of changed agents: correcting an overbroad complexity promise}
\textbf{Problem ID:} \texttt{C78-293}\quad \textbf{JEL:} \texttt{C78}\par
% BEGIN ECONOMICS STATEMENT
\label{problem:OP-0199}
% GAME_THEORY_PROBLEM: {"local_id": "GT-069", "original_number": 69, "kind": "H", "entry_kind": "statement_only", "published": false, "review_required": true, "category": "game_theory"}
\label{game:GT-069}
{\sffamily\small\color{muted}
\noindent\textbf{ID: GT-069.}\quad\textbf{Category:} Robust stable matching complexity.
\quad\textbf{Kind: H.}\quad\textbf{Status:} The catalogue's broad promise
retains NP-complete instances; narrower regimes form a research agenda.
\par}
\subsection*{Definitions and mathematical statement}
Let \(W,F\) have \(n\) members each. Profiles \(A,B\) give strict complete
opposite-side rankings. A matching is a bijection \(\mu:W\to F\);
stability in a profile means there is no worker--firm pair that both
strictly prefer to their current partners. For bounds \(p,q\leq n\),
the \((p,q)\) robust-existence problem asks whether
\(\mathcal M_A\cap\mathcal M_B\ne\varnothing\), assuming that at most
\(p\) worker lists and at most \(q\) firm lists differ.

\paragraph{Correction to the stated promise.}
The family specified only by
\[
 2\leq p,q\leq n,\qquad p+q<2n,\qquad p+q\longrightarrow\infty
\]
is still NP-complete; the last condition is asymptotic, not a property
of one finite input.
Indeed, pad an arbitrary order-\(r\) robust-existence instance by one
worker and one firm. They rank each other first in both profiles and
have identical remaining lists; all old agents place the new participant
last. The new order is \(n=r+1\), and \(p=q=r\) are valid bounds, so
\(p+q=2n-2<2n\).
Every stable matching must pair the new mutually first-choice agents.
Restriction to the old agents, and extension by the new pair, preserve
stability in both profiles. This is a polynomial reduction from the
general problem, which [S1] identifies as NP-complete.
Consequently, polynomial-time solvability of this whole promised family
is equivalent to \(\mathbf P=\mathbf{NP}\).

\paragraph{Remaining formalized agenda (R).}
For specified unbounded functions \(p(n),q(n)\), classify the restricted
robust-existence problems by polynomial-time algorithms or hardness
results. Particular growth restrictions must be stated; merely excluding
\((n,n)\) does not create an unresolved tractable class.

\subsection*{Short English statement}
Leaving one unchanged agent on each side does not remove general hardness.
Which more restrictive growing-change regimes are tractable?

\subsection*{Variants and scope boundaries}
The parameters are upper bounds on changed lists, not necessarily exact
counts. Fixed \(p+q\) and a growing parameter are different regimes.
Membership in NP follows by checking a proposed matching in both profiles.

\subsection*{Source relationship and status}
[S1, Remark 17] motivates the intermediate-regime question.
The padding argument is an editorial correction to the catalogue's
overbroad formulation, not a theorem claimed in that remark.

\subsection*{References}
\begingroup\footnotesize\raggedright
\noindent[S1] Rohith Reddy Gangam, Tung Mai, Nitya Raju, and
Vijay V.~Vazirani,
\emph{Stable Matching: Dealing with Changes in Preferences}, 2023;
arXiv:2304.02590v3, 2025.
\url{https://arxiv.org/html/2304.02590v3}.
Locator: Section 4.4, Remark 17 (general-case NP-completeness,
the XP bound, and the growing-change question).
\par\endgroup

\subsection*{Common conventions: Source mathematical conventions}

\textbf{Finite games.} For a finite set $A$, $\Delta(A)$ is its probability
simplex. Players choose independent mixed strategies in Nash equilibrium;
correlation is allowed only when explicitly part of the solution concept.
In a game with payoff functions $u_i$ and strategy profile $x$, an additive
$\varepsilon$ Nash equilibrium satisfies
\[
 u_i(x)\geq u_i(x_i',x_{-i})-\varepsilon
 \quad\text{for every player $i$ and every admissible deviation $x_i'$}.
\]
For numerical approximation, payoffs are normalized as locally specified.
An $\varepsilon$ payoff residual is distinct from distance at most
$\varepsilon$ to the set of exact equilibria. Point convergence, convergence
of empirical play, and convergence in distance to an equilibrium set are
also distinct.

\textbf{Computational representations.} Unless stated otherwise, a complexity
question takes rational data with binary encoding; polynomial time is measured
in total bit length. A fixed-accuracy polynomial algorithm is not necessarily
polynomial in $1/\varepsilon$, and neither is necessarily polynomial in
$\log(1/\varepsilon)$. Succinct games and value, demand, or sampling oracles
must declare their interfaces. Exact real solutions require an explicit
representation; an unspecified list of arbitrary real numbers is not a
finite computational output. A claimed algorithmic barrier is meaningful
only for its specified model and reductions.

\textbf{Dynamic games.} Histories, information, observation, and allowable
randomization are part of the model. A stationary strategy depends only on
the current state; a positional strategy in a deterministic graph game
chooses one action at each controlled vertex. Nash equilibrium at an initial
state and equilibrium after every history are different requirements.
An expectation of a pathwise $\liminf$ payoff, a $\liminf$ of expected averages,
a limiting expected average, and a uniform finite-horizon equilibrium are
different criteria. None is silently substituted for another.

\textbf{Allocation and mechanisms.} A complete allocation assigns every item
exactly once unless otherwise stated. Zero-valued goods, negative values,
capacity constraints, feasibility, lotteries, and copies of goods affect
fairness definitions and are specified locally. Dominant-strategy incentive
compatibility, Bayesian incentive compatibility, universal truthfulness,
and truthfulness in expectation impose different inequalities. Whenever
an objective ratio has zero denominator, its convention is stated or those
instances are excluded.

\textbf{Infinite-dimensional models.} Expectations, integrals, stochastic
equations, and optimization objectives require their local regularity,
integrability, filtrations, and admissible domains. A backward--forward
mean-field system is a boundary-value problem; it is not an initial-value
semiflow without an additional construction. Long-time, large-population,
and vanishing-noise limits require an order of limits and a convergence
topology. If a minimum need not be attained, the objective is its infimum
or supremum and the approximation requirement is stated separately.
% END ECONOMICS STATEMENT
\end{document}
